% !TEX program = xelatex
\documentclass[aspectratio=169,12pt]{beamer}

% ===== Beamer Theme =====
\usetheme{Madrid}
\usecolortheme{default}

% ===== Chinese Font Support =====
\usepackage{xeCJK}
\setCJKmainfont{Microsoft JhengHei}
\setCJKsansfont{Microsoft JhengHei}
\setCJKmonofont{Microsoft JhengHei}

% ===== Packages =====
\usepackage{graphicx}
\usepackage{booktabs}
\usepackage{tikz}
\usepackage{fontawesome5}

% ===== Colors =====
\definecolor{ntuhblue}{RGB}{0,51,102}
\definecolor{ntuhred}{RGB}{153,0,0}
\definecolor{ethicsgreen}{RGB}{0,102,51}
\definecolor{warningorange}{RGB}{204,102,0}

\setbeamercolor{title}{fg=ntuhblue}
\setbeamercolor{frametitle}{fg=ntuhblue}
\setbeamercolor{structure}{fg=ntuhblue}
\setbeamercolor{block title}{bg=ntuhblue,fg=white}
\setbeamercolor{block body}{bg=ntuhblue!10}
\setbeamercolor{alerted text}{fg=ntuhred}

% ===== Beamer Settings =====
\setbeamertemplate{navigation symbols}{}
\setbeamertemplate{footline}[frame number]
\setbeamerfont{frametitle}{size=\large}

% ===== Title Info =====
\title[醫法倫討論]{醫學倫理法律討論會}
\subtitle{主持人評論 Chair's Commentary}
\author{謝慕揚 MD, PhD, FESC}
\institute{國立臺灣大學醫學院附設醫院新竹臺大分院}
\date{2026年1月16日}

\begin{document}

% ===== Title Slide =====
\begin{frame}
\titlepage
\end{frame}

% ===== Slide 2: Case Recap =====
\begin{frame}{個案回顧 Case Recap}
\textbf{79歲男性 — 關鍵時間線}

\begin{columns}[T]
\begin{column}{0.48\textwidth}
\textcolor{ntuhred}{\textbf{急性期}}
\begin{itemize}
    \item 10/10: STEMI + cardiogenic shock
    \item 緊急PCI（\alert{兩位醫師同意}）
    \item 無身分、無健保、無家屬
\end{itemize}

\vspace{0.3cm}
\textcolor{ethicsgreen}{\textbf{ICU病程}}
\begin{itemize}
    \item 10/14: TB PCR (+) → HERZ
    \item 10/17: IABP removal
    \item 10/29: Extubation 成功
\end{itemize}
\end{column}

\begin{column}{0.48\textwidth}
\textcolor{warningorange}{\textbf{倫理挑戰}}
\begin{itemize}
    \item 10/21: 社工失聯家屬
    \item 10/27: 哥哥回電後拒接
    \item 11/05: \alert{申訴醫療倫理委員會}
\end{itemize}

\vspace{0.3cm}
\textcolor{ntuhblue}{\textbf{結果}}
\begin{itemize}
    \item 11/11: 出院 (MBD)
    \item 總費用: \textbf{\$1,208,467}
    \item 全數由社福資源補助
\end{itemize}
\end{column}
\end{columns}
\end{frame}

% ===== Slide 3: Four Ethical Dilemmas =====
\begin{frame}{四大倫理困境 Four Ethical Dilemmas}

\begin{columns}[T]
\begin{column}{0.48\textwidth}
\begin{block}{1. 自主權 Autonomy}
緊急同意權 — 誰能代為決定？\\
\small{病人昏迷、家屬缺位時的醫療決策權}
\end{block}

\vspace{0.3cm}

\begin{block}{2. 行善 Beneficence}
積極治療 vs. 可能的無效醫療\\
\small{病史不明時，治療強度如何拿捏？}
\end{block}
\end{column}

\begin{column}{0.48\textwidth}
\begin{block}{3. 不傷害 Non-maleficence}
Time-Limited Trial 何時該停止？\\
\small{避免過度醫療造成的痛苦}
\end{block}

\vspace{0.3cm}

\begin{block}{4. 公平正義 Justice}
無健保病人的資源分配\\
\small{\$120萬 — 對其他納稅人公平嗎？}
\end{block}
\end{column}
\end{columns}

\vspace{0.5cm}
\centering
\textit{「無知之幕」(Veil of Ignorance) — 若你不知道自己會是誰...}
\end{frame}

% ===== Slide 4: Commentary - Consent & DNR =====
\begin{frame}{評論：緊急同意與DNR}

\begin{columns}[T]
\begin{column}{0.48\textwidth}
\textcolor{ethicsgreen}{\faCheck\ \textbf{本案處置合理}}
\begin{itemize}
    \item \textbf{醫療法第60條}：危急病人不得拖延
    \item 兩位醫師共同決定
    \item 事後倫理委員會審查
    \item 保留完整書面紀錄
    \item 結果：病人成功脫離險境
\end{itemize}
\end{column}

\begin{column}{0.48\textwidth}
\textcolor{warningorange}{\faExclamationTriangle\ \textbf{值得深思的灰色地帶}}
\begin{itemize}
    \item 「\alert{失聯}」$\neq$「\alert{無親屬}」
    \item 安寧條例第7條第3項適用時機
    \item 「實質無代理人」認定標準
    \item 違反第7條：罰鍰6-30萬或停業
\end{itemize}

\vspace{0.3cm}
\textcolor{ntuhblue}{\small{\textbf{ATS/AGS建議}：盡職搜索 + 書面佐證 + 最大利益標準}}
\end{column}
\end{columns}

\end{frame}

% ===== Slide 5: Commentary - TLT & Resources =====
\begin{frame}{評論：限時嘗試治療與資源分配}

\begin{columns}[T]
\begin{column}{0.48\textwidth}
\textcolor{ethicsgreen}{\textbf{Time-Limited Trial 成功案例}}
\begin{itemize}
    \item 初期病史不明 → 先積極治療
    \item 設定明確觀察期限與目標
    \item 動態評估，團隊共識
    \item \textbf{病情改善} → 驗證TLT價值
\end{itemize}

\vspace{0.3cm}
\textcolor{ntuhred}{\textbf{若持續惡化？}}
\begin{itemize}
    \item 無法「有意義互動」
    \item 不再升級 (No escalation)
    \item 考慮撤除維生裝置
\end{itemize}
\end{column}

\begin{column}{0.48\textwidth}
\textcolor{ntuhblue}{\textbf{資源分配的現實}}

\begin{itemize}
    \item \textbf{行善原則}：提供必要治療
    \item \textbf{公平正義}：健保標準為參考
    \item 提供「基本」而非「最好」？
\end{itemize}

\vspace{0.3cm}
\textcolor{ethicsgreen}{\textbf{本案：多方資源成功介入}}
\begin{itemize}
    \item 台大社福基金：\$40萬
    \item TB補助計畫：\$44萬
    \item 社會局老人保護：\$36萬
\end{itemize}
\end{column}
\end{columns}

\end{frame}

% ===== Slide 6: Discussion Questions 1-3 =====
\begin{frame}{討論問題 Discussion Questions (1/2)}

\begin{alertblock}{Q1: 緊急同意機制}
「兩位醫師共同決定」進行緊急手術是否足夠？\\
是否需要建立更明確的\textbf{院內規範}或\textbf{事後審查機制}？
\end{alertblock}

\vspace{0.3cm}

\begin{alertblock}{Q2: 代理人缺位認定}
家屬「拒接電話」與「完全無親屬」在倫理與法律上有何不同？\\
\textbf{多久}或\textbf{何種情況}可認定為「實質無代理人」？
\end{alertblock}

\vspace{0.3cm}

\begin{alertblock}{Q3: TLT決策點}
本案採用限時嘗試治療，若病人在ICU期間\textbf{持續惡化}，\\
何時應該停止升級治療？判斷標準為何？
\end{alertblock}

\end{frame}

% ===== Slide 7: Discussion Questions 4-5 =====
\begin{frame}{討論問題 Discussion Questions (2/2)}

\begin{alertblock}{Q4: 資源正義}
面對無健保、可能產生\textbf{巨額呆帳}的病人，\\
醫療團隊應提供「基本」還是「最好」的醫療？\\
如何平衡\textbf{行善原則}與\textbf{公平正義}？
\end{alertblock}

\vspace{0.5cm}

\begin{alertblock}{Q5: 系統改善}
本案最終透過多方社福資源解決費用問題。\\
作為醫療體系，我們可以建立什麼\textbf{機制}來更早期識別並處理\\
類似「\textbf{無身分、無健保、無家屬}」的三無個案？
\end{alertblock}

\vspace{0.5cm}
\centering
\textcolor{ntuhblue}{\large{\textit{歡迎各位先進分享看法}}}

\end{frame}

% ===== Slide 8: Take-home Points =====
\begin{frame}{結語 Take-home Points}

\begin{enumerate}
    \item \textbf{緊急情況}：生命優先，但需\alert{完整紀錄}與\alert{事後審查}

    \vspace{0.3cm}

    \item \textbf{代理人缺位}：盡職搜索 + 書面佐證 + 倫理委員會介入

    \vspace{0.3cm}

    \item \textbf{TLT原則}：設定期限、動態評估、團隊共識

    \vspace{0.3cm}

    \item \textbf{資源分配}：健保標準為基準，善用社福資源

    \vspace{0.3cm}

    \item \textbf{預防勝於治療}：建立早期識別機制
\end{enumerate}

\vspace{0.5cm}
\centering
\begin{block}{}
\centering
\textbf{感謝PGY1同學們的精彩報告}\\
\small{陳弘益、張家斌、陳品卉、林鎧伊、黃瀚毅}
\end{block}

\end{frame}

\end{document}
